\section{Content-Typed Endpoints}
\label{sec:content}

Three endpoints look inside a file: inline images, inline PDFs, and archive listings.  Each is
defined by a test on the file's \emph{leading bytes}, never its name, and each opens the file
through the guard, so every property of Section~\ref{sec:guard} applies to it unchanged.  A
fourth, Quick Look pictures, is different in kind: macOS reads the document, not
\texttt{fsb} (Section~\ref{sec:quicklook}).

\subsection{Images}

An image is served inline only if its first bytes are a PNG, JPEG, GIF or WebP signature.
The response is sent with a policy that sandboxes it and allows no resources, so even a
crafted file has nothing to reach.  SVG is never accepted, because it can carry script, and
neither is HTML.  A file named \code{.png} that is not one is refused (415) and the interface
falls back to the ordinary file preview.

\subsection{PDFs}

A PDF is served inline only if it begins with \code{\%PDF-}.  It cannot be sent with the
\code{sandbox} policy that images get, because browsers then refuse to display it.  The
protection is instead the byte check, framing headers that allow only the interface's own
pages to embed it, a policy that forbids script in the response, and the browser's own PDF
viewer, which runs in an isolated component of its own.  The interface asks first, with a
one-byte range request, whether the file really is a PDF, and only then points a frame at it.

A viewer may take the keyboard focus by itself (Chrome's does, for an encrypted PDF's
password field), and once focus is inside the frame no key event reaches the interface.  The
preview pane records whether the pointer is over the frame and whether the last key was Tab;
when the window loses focus to the frame with neither, focus goes back to the list
(\fq{PRV-8}).

\subsection{Quick Look pictures}
\label{sec:quicklook}

For iWork and Office documents the preview is the PNG that Apple's \code{qlmanage -t}
draws, which is what Finder shows.  This is the one place \texttt{fsb} hands a file to
another program, and that program opens what it is given by itself, resolving more than the
guard can see: a Finder alias is a small regular file, marked by its \code{FinderInfo}
attribute, that Quick Look follows to its target.  An independent review showed that an
alias named \code{report.docx} made Quick Look draw a file in a denied folder, and one
outside the roots.  Quick Look is therefore never given the user's file.

\begin{enumerate}
  \item \textbf{A private copy.}  \code{guard.CopyForQuickLook} opens the document
        exactly as every endpoint does (roots, deny rules before and after, hard links to
        denied files), requires an iWork or Office extension on the real path, refuses a
        cloud-only document, and copies the bytes, from the descriptor it verified, into a
        fresh temporary folder that only the user can read.  A copy has no extended
        attributes and no resource fork, so an alias's copy is just its bytes, which Quick
        Look draws, at most, as a blank page; and nothing can be swapped between the check
        and the drawing.
  \item \textbf{Packages.}  A package document is walked without following links.  Every
        file in it is opened through the guard by its full path, and must still lie inside
        the package (a folder swapped for a link after the walk listed it leads elsewhere,
        and is caught here).  A part that is denied, a link, a hard link to a denied file or
        not a regular file is left out of the copy, as a listing leaves it out, so the
        answer does not reveal it (the same review found that refusing the whole package
        did); a cloud-only part refuses the package.
  \item \textbf{Bounds.}  A copy is limited to 256~MB and 5{,}000 entries.  Copying and
        drawing share two slots.  \code{qlmanage} never finishes on a document it cannot
        draw, so it runs under a 5-second deadline and is killed when the browser abandons
        the request.  Exactly one PNG under 20~MB must come out, or there is no picture.
        The picture is served like an image (sandbox policy, \code{nosniff}) and never
        cached, and the temporary folder is removed.
\end{enumerate}

Choosing by extension rather than by bytes (\fq{PRV-2}) is deliberate: Quick Look decides
what the document is, and a package is a folder with no bytes of its own.  The extension only
decides whether to ask.  Keynote's \code{.key} stays covered by the core rule \code{*.key}.

\subsection{Archives}

An archive listing returns names and sizes only, never member data, and nothing is extracted.
The format is decided from bytes: the zip signature, the \code{ustar} marker at offset 257 of a
tar, or a gzip stream that contains a tar.  The limits of \fq{LIM-1} bound a hostile file:

\begin{itemize}
  \item A zip's directory is held in memory by the standard library, so the guard first reads
        the end-of-central-directory record and refuses a zip that claims more entries than it
        will index (and, when the count is stored elsewhere in a very large zip, refuses by
        size).
  \item A tar is streamed, stopping at an entry limit; a gzip-compressed tar is read through a
        reader capped at a number of decompressed bytes, so a bomb costs a bounded amount of
        time and memory, and the listing says it was cut short.
  \item Member names are untrusted: control characters and invalid UTF-8 are replaced, and the
        interface additionally replaces bidirectional characters, and names are only ever
        displayed, never used as paths.
\end{itemize}

\subsection{Rendered Markdown}

Rendering a Markdown file must not run code, load anything remote, or reach the interface.
It is done in three layers, each assuming the one before failed.

\begin{enumerate}
  \item \textbf{Escaping.}  The Markdown is converted by a library configured so that raw HTML
        in the source becomes text; links carry no address at all (an attribute says what
        they mean: a file on the disk, or a web address); images carry no source (a marker names
        a local file); a remote or unsupported image becomes a text placeholder.  A link is
        resolved against the Markdown file's folder to a normalized absolute path, or dropped
        to plain text if it is not a relative path or an \code{http} or \code{https} address.
  \item \textbf{Sanitizing.}  The result goes through DOMPurify with a fixed list of allowed
        tags and attributes.
  \item \textbf{A sandboxed frame.}  The sanitized markup is sent by \code{postMessage} to a fixed
        page, framed with \code{sandbox="allow-scripts"} (so it has an opaque origin and no access
        to the session or the interface).  The page's own policy is \code{default-src 'none'}
        with images limited to \code{data:}, and admits its one script and one style
        \emph{by hash}, computed from the page at start so that a change to the page cannot
        silently invalidate the policy.  The page accepts messages only from its parent.
\end{enumerate}

Links are not followed inside the frame.  A click sends a message to the interface, which
validates it and either changes the location to the target file's folder with the file selected,
or opens a web address in a new tab without a referrer.  Images inside the roots are fetched by
the interface, through the guarded image endpoint, and passed in as \code{data:} URLs, because the
cookie's scope and \code{SameSite} would otherwise stop the opaque-origin frame from fetching them.
Opening a Markdown file therefore makes no request to anything outside the machine.
